\documentclass{article}

% Official NeurIPS 2026 style, unmodified, in template/ (from
% https://media.neurips.cc/Conferences/NeurIPS2026/Formatting_Instructions_For_NeurIPS_2026.zip).
% Build: latexmk -r template/latexmkrc
% [preprint] shows the author block and the acknowledgements, which carry the
% required AI use statement; the default submission mode would hide both.
\usepackage[preprint]{neurips_2026}

% Packages from the official neurips_2026.tex template.
\usepackage[utf8]{inputenc} % allow utf-8 input
\usepackage[T1]{fontenc}    % use 8-bit T1 fonts
\usepackage{hyperref}       % hyperlinks
\usepackage{url}            % simple URL typesetting
\usepackage{booktabs}       % professional-quality tables
\usepackage{amsfonts}       % blackboard math symbols
\usepackage{nicefrac}       % compact symbols for 1/2, etc.
\usepackage{microtype}      % microtypography
\usepackage{xcolor}         % colors

% Added for this paper.
\usepackage{graphicx}       % figures; set widths as a multiple of \linewidth

\title{TITLE}

\author{%
  AUTHOR \\
  AFFILIATION \\
  \texttt{EMAIL} \\
}

\begin{document}

\maketitle

\begin{abstract}
  ABSTRACT
\end{abstract}

% Main text goes here (6-page limit; acknowledgements, references and
% appendices do not count).

\begin{ack}
  % AI use statement (required by the course): name the assistants used and
  % what they were used for.
\end{ack}

% References. natbib is loaded by neurips_2026.sty; add entries to
% references.bib and cite with \citet / \citep.
\bibliographystyle{plainnat}
\bibliography{references}

% \appendix
% \section{...}

\end{document}
